\documentclass[11pt]{amsart}
\usepackage{amsmath,amssymb,amsthm,enumitem,booktabs}
\usepackage[margin=1.05in]{geometry}
\newtheorem{theorem}{Theorem}[section]
\newtheorem{lemma}[theorem]{Lemma}
\newtheorem{proposition}[theorem]{Proposition}
\newtheorem{corollary}[theorem]{Corollary}
\newtheorem{remark}[theorem]{Remark}
\newtheorem{convention}[theorem]{Convention}
\newcommand{\eps}{\varepsilon}
\newcommand{\eu}{\mathrm{e}}
\newcommand{\D}{\mathcal D}
\newcommand{\Def}{\mathrm{Def}}
\newcommand{\pos}[1]{{(#1)_+}}
\newcommand{\kxi}{\kappa_\xi}
\newcommand{\kq}{\kappa_q}
\newcommand{\kmx}{\kappa_{\max}}
\newcommand{\khi}{\kappa_{+}}
\newcommand{\lbar}{\bar\ell}
\newcommand{\RII}{[\textrm{R2}]}
\newcommand{\ZB}{[\textrm{ZB}]}
\newcommand{\ZC}{[\textrm{ZC}]}
\DeclareMathOperator{\Tr}{Tr}

\title[Region II for the small odd cycles $m\in\{9,11,13\}$]%
{The Region-II scalar inequality for the small odd cycles $m\in\{9,11,13\}$\\
(final assembled section)}
\author{}
\date{2026-07-17}

\begin{document}
\maketitle

\begin{abstract}
For each odd cycle length $m\in\{9,11,13\}$ we prove the Region-II scalar
inequality $R_m\le C_m\psi(\xi,\rho)$ on the whole admissible domain, and hence
(via the $m$-uniform reduction of $\RII$) the odd-cycle homomorphism bound
$t(C_m,W)\ge p^m-pq^{m-1}$ in Region~II.  The proof partitions the admissible
domain into three analytic zones---small~$e$ (Zone~A), the $\xi\ge1$ moderate
strip (Zone~B), and the $\xi\le1$ slice including the Tur\'an corner
(Zone~C)---each closed by a fully displayed, fixed (non\-adaptive) certificate.
Every displayed constant is an exact rational rounded in a stated safe
direction; the companion script \texttt{validate\_m91113.py} re-verifies all of
them numerically and in exact rational arithmetic, and confirms that the three
zones tile the domain.
\end{abstract}

\section{Statement, notation and result status}\label{sec:intro}

We work in the setting of the Region-II manuscript
\texttt{paper\_region2\_v2.tex} (cited as $\RII$), whose $m\ge15$ analysis is
carried out by the two verified analytic zone notes
\texttt{zoneB\_analytic.tex} ($\ZB$) and
\texttt{proof\_zoneC\_analytic\_v2.tex} ($\ZC$); the present note specialises
those templates to the three fixed short lengths $m\in\{9,11,13\}$.  All
notation is that of $\RII$: $q\in(\tfrac13,\tfrac12)$, $p=1-q$,
$\alpha\in(q,r(q)]$ the frontier eigenvalue with
$r(q)=\tfrac{\sqrt{q^2+4q}-q}2$, $L=\sqrt{pq-\alpha^2}$, $d=\alpha-q$,
$e=1-2\alpha$, $f=\alpha-L$, $\kappa=d/e$, $x=\alpha/p$, $\tau=q/\alpha$,
$s=q/p$, $y=L/p$, $\ell=L/\alpha$, $u=1-x$, $\delta=\alpha-\tfrac13$, and
\begin{equation}\label{eq:defs}
 k_m(\lambda)=\frac{p^{m-1}-\lambda^{m-1}}{p+\lambda},\quad
 A_m=2L^{m-2}+m\,k_m(\alpha),\quad
 B_m=2L^{m-2}+m\,k_m(L),\quad
 R_m=\alpha^m+L^m-pq^{m-1}.
\end{equation}
We use the $(e,\kappa)$ chart of $\RII$ (Lem.~\texttt{chart}): the admissible
$(q,\alpha)$ are exactly
\begin{equation}\label{eq:chart}
 0<e<\tfrac13,\qquad 0<\kappa\le\kmx(e)=\tfrac{1-e}{1+e},\qquad
 \kappa<\kq(e)=\tfrac{1-3e}{6e},
\end{equation}
and on it $\alpha=\tfrac{1-e}2$, $L^2=\alpha e-d(d+e)$,
$\xi=\tfrac{(1-e)^2\kappa}{e}=\tfrac{4\alpha^2 d}{e^2}$,
$\delta=\tfrac{1-3e}6$.  The Huber payment of $\RII$ (Def.~\texttt{huber}) is
\begin{equation}\label{eq:target}
 R_m\ \le\ C_m\,\psi(\xi,\rho),\qquad
 C_m=\frac{\sqrt{2\alpha}}{4\alpha^2}\,B_m f\,e^2,\quad
 \rho=\frac{A_m}{B_m}\,\frac{\sqrt\alpha}{2\sqrt2\,f},\quad
 \psi=\max_{0\le\lambda\le1}\Bigl(\lambda\xi-\tfrac{\lambda^2}{4(\rho+\lambda)}\Bigr),
\end{equation}
and \eqref{eq:target} is the inequality this note establishes for
$m\in\{9,11,13\}$.

\begin{theorem}\label{thm:main}
For every $m\in\{9,11,13\}$ and every admissible $(q,\alpha)$,
inequality \eqref{eq:target} holds.  Consequently
(Remark~\ref{rem:muniform}) the odd-cycle bound $t(C_m,W)\ge p^m-pq^{m-1}$
holds in Region~II for these $m$.
\end{theorem}

\begin{remark}[$m$-uniformity of the $\RII$ reduction]\label{rem:muniform}
The statement of $\RII$ Thm.~\texttt{huber} carries the standing hypothesis
``odd $m\ge15$'', but its proof---and the proofs of the whole reduction chain
it rests on ($\RII$ Prop.~\texttt{shift}, Lem.~\texttt{no-frontier},
Lem.~\texttt{one-frontier}, Lem.~\texttt{forced-variance},
Prop.~\texttt{master-defect}, Prop.~\texttt{psi-forms})---is $m$-pointwise: $m$
enters only through the quantities $A_m,B_m,R_m$ of \eqref{eq:defs}, and the
only structural inputs are $0<A_m<B_m$ ($\RII$ Rem.~\texttt{knegative}) and the
frontier geometry $L<q<\alpha<p$, valid for every odd $m\ge3$ on the admissible
domain.  ($A_m<B_m$ since $k_m$ is strictly decreasing on $[0,p]$ and
$L<\alpha$; $A_m>0$ since $\alpha<p$ gives $k_m(\alpha)>0$.)  Hence the
implication ``\eqref{eq:target} on the admissible domain $\Rightarrow$
$t(C_m,W)\ge p^m-pq^{m-1}$ in Region~II'' holds verbatim for every odd
$m\ge3$, and Theorem~\ref{thm:main} may invoke it for $m\in\{9,11,13\}$.
\end{remark}

\emph{Status of the proof.}  Every step below is analytic in the sense fixed by
$\RII$/$\ZB$/$\ZC$: the certificate relaxation and positivity gate
(\S\ref{sec:relax}); the small-$e$ range $e\le\tfrac1{60}$
(Proposition~\ref{prop:smalle}, Zone~A); the $\xi\ge1$ moderate strip via a
$26$-row displayed table over a \emph{fixed} rational cover
(\S\ref{sec:zoneB}, Zone~B); and the $\xi\le1$ slice, including the Tur\'an
corner, via displayed $\varphi$-row and dual-certificate tables over a
\emph{fixed} cover (\S\ref{sec:zoneC}, Zone~C).  No adaptive/greedy interval
subdivision appears anywhere.  Throughout, a terminating decimal denotes that
exact rational; every decimal comparison is rounded in the safe direction with
the underlying exact comparison stated, and all displayed inequalities are
additionally machine-verified by \texttt{validate\_m91113.py}.

\begin{convention}[safe decimals]\label{conv:safe}
A quantity that must be bounded above is written as an exact rational rounded
up (marked by ``$\le$''); one bounded below is rounded down.  Each displayed
row is a finite exact rational computation certifying its inequality; the
script re-verifies every row from the unrounded rationals.
\end{convention}

\section{The relaxed certificate and the positivity gate}\label{sec:relax}

\begin{lemma}[Relaxed $\lambda=1$ certificate]\label{lem:relax}
On the admissible domain, with $w:=\sqrt{2\alpha}$,
$\rho=\dfrac{A_m w}{4B_m f}$, and consequently
\begin{equation}\label{eq:Tm}
 C_m\psi(\xi,\rho)\ \ge\ \sqrt{2\alpha}\,B_mf\Bigl(d-\frac{e^2}{16\alpha^2(1+\rho)}\Bigr)
 \ \ge\ \sqrt{2\alpha}\,B_mfd-\frac{e^2B_m^2f^2}{4\alpha^2A_m}.
\end{equation}
Hence \eqref{eq:target} follows from
\begin{equation}\label{eq:Tm2}
 (T_m):\qquad 4\alpha^2A_m R_m\ \le\ 4\alpha^2 A_m\sqrt{2\alpha}\,B_mfd-e^2B_m^2f^2 .
\end{equation}
\end{lemma}

\begin{proof}
By $\RII$ Def.~\texttt{huber}, $\rho=\tfrac{A_m}{B_m}\tfrac{\sqrt\alpha}{2\sqrt2 f}$, and
$\tfrac{\sqrt\alpha}{2\sqrt2}=\tfrac{\sqrt{2\alpha}}4$, giving $\rho=A_mw/(4B_mf)$.
The first inequality in \eqref{eq:Tm} is the dual certificate $\lambda=1$
($\RII$ eq.~\texttt{cert-II}).  For the second, $1+\rho\ge\rho$ gives
\[
 \sqrt{2\alpha}\,B_mf\cdot\frac{e^2}{16\alpha^2(1+\rho)}
 \le \sqrt{2\alpha}\,B_mf\cdot\frac{e^2}{16\alpha^2}\cdot\frac{4B_mf}{A_m\sqrt{2\alpha}}
 =\frac{e^2B_m^2f^2}{4\alpha^2A_m}. \qedhere
\]
\end{proof}

\begin{lemma}[Positivity gate; all odd $m\ge3$]\label{lem:gate}
If $R_m>0$ then
\begin{equation}\label{eq:gate}
 (m-1)\kappa\ >\ x\bigl(1+\kappa-\ell^{m-2}\bigr).
\end{equation}
\end{lemma}

\begin{proof}
By $\RII$ eq.~\texttt{defect-form}, $R_m=\alpha^{m-1}(\alpha-p\tau^{m-1})+L^m$, so
$R_m>0$ means $\tau^{m-1}<x(1+\ell^m)$, since $L^m/(p\alpha^{m-1})=x\ell^m$.
Bernoulli's inequality gives $\tau^{m-1}=(1-d/\alpha)^{m-1}\ge1-(m-1)d/\alpha$,
whence $(m-1)\tfrac d\alpha>1-x-x\ell^m$.  Multiply by $\alpha/e$ and use
$1-x=u=\tfrac{d+e}p=\tfrac{x(d+e)}\alpha$, so that $(1-x)\alpha/e=x(1+\kappa)$;
finally $\ell^m\alpha/e\le\ell^{m-2}$ because $\ell^2\le e/\alpha$
($\RII$ Lem.~\texttt{elementary}).
\end{proof}

\begin{corollary}[Gate on the small-$e$ range]\label{cor:gate-smalle}
Let $3\le m\le13$ be odd, $e\le1/60$, and $R_m>0$.  Then, with the standing
bounds $x\ge\tfrac{1-e}{1+3e}\ge\tfrac{59}{63}>0.9365$ and
$\ell\le\sqrt{2e/(1-e)}\le1.4262\sqrt e\le0.1842$,
\begin{equation}\label{eq:gate-smalle}
 \kappa\ >\ \frac{x(1-\ell)}{m-1-x}\ \ge\ \frac{0.763}{m-1}\ \ge\ 0.0635,
 \qquad
 \xi=\frac{(1-e)^2\kappa}e\ \ge\ 3.68,
\end{equation}
and for $m=3$ moreover $\kappa>0.7183$.  In particular the region
$\{\,e\le1/60,\ \xi\le1,\ R_m>0\,\}$ is empty for $3\le m\le13$, and
\begin{equation}\label{eq:eps-smalle}
 \eps:=\frac{e^2}{16\alpha^2(1+\rho)d}=\frac{e}{4(1-e)^2(1+\rho)\kappa}
 \ \le\ \min\Bigl(0.068,\ \frac{e}{4(1-e)^2\kappa}\Bigr).
\end{equation}
\end{corollary}

\begin{proof}
Since $m\ge3$ and $\ell<1$, $\ell^{m-2}\le\ell$, so \eqref{eq:gate} gives
$\kappa(m-1-x)>x(1-\ell)$, i.e.\ the first bound; then
$x(1-\ell)\ge0.9365\cdot0.8158>0.763$ and $m-1-x<m-1$.  For $m=3$,
$t\mapsto t/(2-t)$ is increasing, so
$\kappa>0.9365\cdot0.8158/(2-0.9365)>0.7183$.  Next,
$\xi\ge(59/60)^2\cdot\tfrac{0.763}{12}\cdot60>3.68$ using $m\le13$; in
particular $\xi>1$.  Finally $\eps=1/(4\xi(1+\rho))\le1/(4\cdot3.68)<0.068$ and
$\rho\ge0$; the equality in \eqref{eq:eps-smalle} is the chart identity
$\xi=(1-e)^2\kappa/e$.
(Decimal replacements: $0.9365\cdot0.8158=0.76399\ldots>0.763$,
$0.9365\cdot0.8158/1.0635=0.71837\ldots>0.7183$,
$(59/60)^2\cdot0.763\cdot5=3.68889\ldots>3.68$, $1/14.72=0.06793\ldots<0.068$.)
\end{proof}

\section{Zone A: the small-$e$ range $e\le1/60$}\label{sec:zoneA}

\begin{proposition}[Zone A]\label{prop:smalle}
Let $3\le m\le13$ be odd, $e\le1/60$, $R_m>0$.  Then
$R_m\le\sqrt{2\alpha}B_mf\bigl(d-\frac{e^2}{16\alpha^2(1+\rho)}\bigr)\le C_m\psi$,
so \eqref{eq:target} holds.  (If $R_m\le0$, \eqref{eq:target} is trivial.)
\end{proposition}

\begin{lemma}[Reduction; $m$-adapted $\RII$ Lem.~\texttt{reduceA}]\label{lem:reduce}
Under the hypotheses of Proposition~\ref{prop:smalle}, the first inequality
follows from
\begin{equation}\label{eq:red}
 \sqrt{1-e}\,\frac{(1-\ell)(1-y^{m-1})(1-\eps)}{1+y}
 \ \ge\
 x^{m-2}\Bigl[\frac{m-1}{mx}-\frac{1+1/\kappa}{m}\Bigr]+\Lambda,
 \qquad
 \Lambda:=\frac{x^{m-2}(e/\alpha)^{m/2-1}}{m\kappa}.
\end{equation}
\end{lemma}

\begin{proof}
Exactly as in $\RII$ Lem.~\texttt{reduceA}:
$B_m\ge mk_m(L)=\frac{mp^{m-2}(1-y^{m-1})}{1+y}$, $\sqrt{2\alpha}=\sqrt{1-e}$,
$f=\alpha(1-\ell)$ give payment $\ge md\alpha p^{m-2}\cdot\mathrm{LHS}\eqref{eq:red}$;
Bernoulli $\tau^{m-1}\ge1-(m-1)d/\alpha$ and $L^m\le(\alpha e)^{m/2}$ give
$R_m\le md\alpha p^{m-2}\cdot\mathrm{RHS}\eqref{eq:red}$.
(No $m\ge15$ constant enters.)
\end{proof}

\begin{proof}[Proof of Proposition~\ref{prop:smalle}]
Standing bounds for $e\le1/60$, each with its exact justification:
$x\ge\tfrac{1-e}{1+3e}\ge\tfrac{59}{63}>0.9365$ (from $d\le e$);
$\ell^2\le e/\alpha=\tfrac{2e}{1-e}\le\tfrac{120}{59}e\le2.04e$, hence
$\ell\le1.4262\sqrt e\le0.1842$
($1.4262^2=2.03404\ldots>120/59=2.03389\ldots$; $1.4262/\sqrt{60}<0.1842$);
$y\le\ell$ (as $p\ge\alpha$); $\sqrt{1-e}\ge1-0.51e$
(as $(1-0.51e)^2-(1-e)=-0.02e+0.2601e^2\le0$ for $e\le0.02/0.2601$);
$1/x=1+(1+\kappa)e/\alpha\le1+2.04(1+\kappa)e$ (as $1/\alpha\le120/59\le2.04$);
$u=1-x=\tfrac{(1+\kappa)e}{p}\ge1.9047(1+\kappa)e$
(as $p\le\tfrac{1+3e}2\le\tfrac{21}{40}$ and $40/21=1.90476\ldots\ge1.9047$);
and Corollary~\ref{cor:gate-smalle}.

\emph{Case $m=3$.}  Using
$\frac{(1-\ell)(1-y^2)(1-\eps)}{1+y}\ge(1-\ell)(1-y)(1-y^2)(1-\eps)$ and
$(1-a_1)\cdots(1-a_k)\ge1-\sum a_i$ on $[0,1]$,
\[
 \mathrm{LHS}\eqref{eq:red}\ \ge\ 1-2\ell-y^2-\eps-0.51e
 \ \ge\ 1-0.3684-0.034-0.068-0.0085\ >\ 0.52,
\]
where $y^2\le\ell^2\le2.04e\le2.04/60=0.034$.  For $m=3$ the bracket is
$\tfrac{2}{3x}-\tfrac{1+1/\kappa}3$, so
$\mathrm{RHS}\eqref{eq:red}=\tfrac23-\tfrac{x(1+1/\kappa)}3+\Lambda$ with
$\Lambda=\tfrac{x\sqrt{e/\alpha}}{3\kappa}$.  Using $\kappa>0.7183$,
$1+1/\kappa>2$, $x\ge0.9365$, $x\le1$, $\sqrt{e/\alpha}\le1.4262\sqrt e$ and
$1.4262/(3\cdot0.7183)=0.66183\ldots\le0.662$:
\[
 \mathrm{RHS}\eqref{eq:red}\le\tfrac23-\tfrac{0.9365\cdot2}{3}+0.662\sqrt e
 \le\tfrac23-0.6243+0.662\cdot0.1291<0.128<0.52,
\]
($\sqrt{1/60}<0.1291$ as $0.1291^2=0.01666\,681>1/60$).  So \eqref{eq:red}
holds for $m=3$.

\emph{Case $5\le m\le13$.}  Put $z:=(m-2)u$, $c_m:=\tfrac{m-3}{2(m-2)}$,
$\beta_m:=\tfrac{1.0711}{m-2}$ (so $2.04(1+\kappa)e\le\beta_m z$, since
$\beta_mz\ge1.0711\cdot1.9047(1+\kappa)e=2.04012\ldots(1+\kappa)e$), and
$\Gamma_0(e):=2.8524\sqrt e+0.51e+(2.04e)^2$,
$\sigma_m(e):=1.21\,(2.04e)^{(m-2)/2}$.  Then
$2\ell+0.51e+y^{m-1}\le\Gamma_0$ (using $y^{m-1}\le\ell^4\le(2.04e)^2$ for
$m\ge5$), and $\Lambda\le\sigma_mx^{m-2}$: indeed
$(e/\alpha)^{(m-2)/2}\le(2.04e)^{(m-2)/2}$, and by
Corollary~\ref{cor:gate-smalle} and $m\le13$,
$\tfrac1{m\kappa}<\tfrac{m-1}{0.763\,m}\le\tfrac{12}{0.763\cdot13}=1.20979\ldots<1.21$.
Since $x^{m-2}\le(1+z+c_mz^2)^{-1}$ ($\RII$ eq.~\texttt{E-a}) and
$1/x\le1+\beta_mz$, $\tfrac1{mx}\ge\tfrac1m$, \eqref{eq:red} follows from
\begin{equation}\label{eq:F}
 F:=\bigl(1-\Gamma_0-\eps\bigr)(1+z+c_mz^2)-\tfrac1x+\tfrac{2+1/\kappa}m-\sigma_m\ \ge\ 0 .
\end{equation}
We record two endpoint evaluations, rounded up
($2.8524/\sqrt{60}<2.8524/7.7459<0.36825$, since $7.7459^2=59.99896681<60$;
$0.51/60=0.0085$; $(2.04/60)^2=0.001156$; $1.293/60=0.02155$):
\begin{equation}\label{eq:G-ends}
 \Gamma_0(e)\le\Gamma_0(1/60)<0.378,\qquad
 \Gamma_1(e):=\Gamma_0(e)+1.293e\le\Gamma_1(1/60)<0.3995,
\end{equation}
both $\Gamma_0,\Gamma_1$ increasing in $e$; and
$\sigma_m(e)\le\sigma_5(1/60)=1.21\cdot(0.034)^{3/2}<0.0076$ for all
$5\le m\le13$.

\emph{Sub-case $\kappa<\tfrac15$.}  Here $\tfrac{2+1/\kappa}m\ge\tfrac7m$, and
by \eqref{eq:eps-smalle} $\eps\le0.068$.  Since
$1-\Gamma_0-0.068-\beta_m\ge1-0.378-0.068-0.35704>0.19>0$
($\beta_m\le\beta_5=1.0711/3=0.35703\ldots$), dropping $c_mz^2\ge0$,
\[
 F\ \ge\ z\bigl(1-\Gamma_0-0.068-\beta_m\bigr)+\tfrac7m-\Gamma_0-0.068-\sigma_m
 \ \ge\ \tfrac7{13}-0.378-0.068-0.0076\ >\ 0.084\ >\ 0 .
\]

\emph{Sub-case $\kappa\ge\tfrac15$.}  Here \eqref{eq:eps-smalle} gives
$\eps\le\tfrac{5e}{4(1-e)^2}\le\tfrac{4500}{3481}e\le1.293e$, so
$1-\Gamma_0-\eps\ge1-\Gamma_1>0$ by \eqref{eq:G-ends}.  Dropping $c_mz^2\ge0$
and using $1/x\le1+\beta_mz$, $\tfrac{2+1/\kappa}m\ge\tfrac3m$ ($\kappa<1$):
$F\ge z(1-\Gamma_1-\beta_m)+\tfrac3m-\Gamma_1-\sigma_m$.  The bracket is
$\ge1-0.3995-0.35704>0.24$, and $z\ge2.2856(m-2)e$ (from $1+\kappa\ge1.2$), so
\begin{equation}\label{eq:Gm}
 F\ \ge\ G_m(e):=2.2856(m-2)e\bigl(1-\Gamma_1(e)-\beta_m\bigr)
 +\tfrac3m-\Gamma_1(e)-\sigma_m(e) .
\end{equation}
Bounding the bracket once more by $\Gamma_1\le0.3995$ and expanding
$\Gamma_1(e)=2.8524\sqrt e+1.803e+4.1616e^2$,
\begin{equation}\label{eq:Ghat}
 G_m(e)\ \ge\ \widehat G_m(e):=\tfrac3m+\bigl(c_{1,m}e-2.8524\sqrt e\bigr)
 -4.1616e^2-\sigma_m(e),
\end{equation}
\begin{equation}\label{eq:c1}
 c_{1,m}:=2.2856(m-2)\bigl(0.6005-\beta_m\bigr)-1.803
 \ =\ 1.3725028\,(m-2)-4.25110616,
\end{equation}
an exact terminating decimal for each $m$ (Table~\ref{tab:smalle}).

\emph{$\widehat G_m$ is strictly decreasing on $(0,1/60]$.}  The terms
$-4.1616e^2$, $-\sigma_m(e)$ and the constant are nonincreasing in $e$; for the
remaining term, for $0<e\le1/60$,
\[
 \frac{d}{de}\bigl(c_{1,m}e-2.8524\sqrt e\bigr)
 =c_{1,m}-\frac{1.4262}{\sqrt e}\ \le\ c_{1,m}-1.4262\sqrt{60}
 \ <\ 10.847-11.045\ <\ 0,
\]
because $c_{1,m}$ increases in $m$ with $c_{1,13}=10.84642464<10.847$, and
$1.4262\sqrt{60}>1.4262\cdot7.745=11.045919>11.045$ ($7.745^2=59.985025<60$).
Hence $\widehat G_m(e)\ge\widehat G_m(1/60)$, and with the safe roundings above
($\sigma_m(1/60)\le\bar\sigma_m$, $c_{1,m}\ge\underline c_{1,m}$ of
Table~\ref{tab:smalle}),
\[
 \widehat G_m(1/60)\ \ge\ \tfrac3m+\tfrac{\underline c_{1,m}}{60}
 -0.36825-0.0011561-\bar\sigma_m\ \ge\ \widehat G_m^{\,\flat}
\]
with the values $\widehat G_m^{\,\flat}$ of Table~\ref{tab:smalle}, all
positive (minimum $0.0383$ at $m=11$).  This proves $F\ge0$ in both sub-cases,
hence \eqref{eq:red}, hence the proposition via Lemma~\ref{lem:reduce} and
Lemma~\ref{lem:relax}.
\end{proof}

\begin{table}[h]\centering
\begin{tabular}{lccccc}\toprule
$m$&5&7&9&11&13\\\midrule
$\beta_m$&$0.35704$&$0.21422$&$0.15302$&$0.11902$&$0.09738$\\
$\underline c_{1,m}$ &$-0.13360$&$2.61140$&$5.35641$&$8.10141$&$10.84642$\\
$\bar\sigma_m$&$0.0076$&$0.000258$&$0.0000088$&$0.0000003$&$0.00000002$\\
$\widehat G_m^{\,\flat}$&$0.2207$&$0.1024$&$0.0531$&$0.0383$&$0.0421$\\
\bottomrule
\end{tabular}
\caption{Small-$e$ constants ($\beta_m$ rounded up; $\underline c_{1,m}$ the
exact value \eqref{eq:c1} truncated toward $-\infty$ at the fifth decimal;
$\bar\sigma_m$ an upper bound for $\sigma_m(1/60)$;
$\widehat G_m^{\,\flat}$ the resulting lower bound for $\widehat G_m(1/60)$,
truncated down).}\label{tab:smalle}
\end{table}

\section{Zone B: the $\xi\ge1$ moderate strip}\label{sec:zoneB}

This section proves \eqref{eq:target} on the $\xi\ge1$ part of Region~II with
$e\ge\tfrac1{60}$, i.e.\ on the strip
\[
 e\in\Bigl[\tfrac1{60},\ e_1\Bigr],\ \ e_1:=\tfrac{2033}{10000},\qquad
 \kappa\in\bigl[\kxi(e),\ \khi(e)\bigr],\quad \khi:=\min(\kmx,\kq),\quad
 \kxi(e):=\tfrac{e}{(1-e)^2},
\]
by specialising $\ZB$ to a single fixed $m$.  ($\xi\ge1$ is exactly
$\kappa\ge\kxi$; the strip is empty for $e>e_1$, where $\kxi>\khi$, so nothing
is claimed there.)  The one structural simplification over $\ZB$ is that, for a
\emph{fixed} $m$, there is no maximisation over cycle length, so $\RII$
Lem.~\texttt{lem:Kmax} (the $m$-maximisation, the envelope
$K^\star=e^{-\Phi-2w(\Phi)}$ and its gate) is \emph{not needed}.  From the chart
Lem.~\texttt{chart} of $\RII$, with $D:=1+e+2\kappa e$,
\begin{equation}\label{zb9:chart}
 x=\frac{1-e}{D},\qquad u=1-x=\frac{2e(1+\kappa)}{D},\qquad
 A:=\frac{x(1+\kappa)}{\kappa}=\frac{(1-e)(1+\kappa)}{\kappa D},
\end{equation}
and $\kmx=\tfrac{1-e}{1+e}$, $\kq=\tfrac{1-3e}{6e}$.  Decimals obey
Convention~\ref{conv:safe}.

\begin{lemma}[reduction to the reduced inequality; fixed $m$]\label{zb9:red}
Let $m\ge3$ be odd and $(q,\alpha)$ admissible with $R_m>0$.  With
$\eps=\tfrac{e^2}{16\alpha^2(1+\rho)d}=\tfrac{e}{4(1-e)^2(1+\rho)\kappa}$,
\eqref{eq:target} holds as soon as
\begin{equation}\label{zb9:reduced}
 \Pi\ :=\ \sqrt{1-e}\,\frac{(1-\ell)(1-y^{m-1})(1-\eps)}{1+y}
 \ \ge\ \frac{x^{m-3}\,(m-1-A)_+}{m}\ +\ \Lambda,
 \qquad
 \Lambda:=\frac{x^{m-2}(e/\alpha)^{m/2-1}}{m\kappa}.
\end{equation}
\end{lemma}

\begin{proof}
This is the $m$-pointwise reduction of $\RII$: the dual certificate
\texttt{cert-II} ($\lambda=1$) gives
$R_m\le\sqrt{2\alpha}B_mf\bigl(d-\tfrac{e^2}{16\alpha^2(1+\rho)}\bigr)$, and the
payment/defect estimates of $\RII$ Lem.~\texttt{reduceA} (used verbatim as
Lemma~\ref{lem:reduce}, ``no $m\ge15$ constant enters'') turn this into
\eqref{zb9:reduced}: one bounds $B_m\ge mk_m(L)=\tfrac{mp^{m-2}(1-y^{m-1})}{1+y}$,
$\sqrt{2\alpha}=\sqrt{1-e}$, $f=\alpha(1-\ell)$ for the payment, and uses
$\tau^{m-1}\ge1-(m-1)d/\alpha$ (Bernoulli) and $L^m\le(\alpha e)^{m/2}$ for the
defect.  The right side is exactly
$x^{m-2}\bigl[\tfrac{m-1}{mx}-\tfrac{1+1/\kappa}{m}\bigr]+\Lambda$, and since
$\tfrac{1+1/\kappa}{m}=\tfrac{A}{mx}$ by \eqref{zb9:chart},
$x^{m-2}\bigl[\tfrac{m-1}{mx}-\tfrac{A}{mx}\bigr]=\tfrac{x^{m-3}(m-1-A)}{m}$,
replaced by its positive part (a larger quantity).
\end{proof}

\begin{remark}[no $m$-maximisation]\label{zb9:rmk:noKmax}
For fixed $m$ the defect is the single value $K(m-1)/x^2$ with
$K(n)=x^n(n-A)/(n+1)$; there is no maximisation over $n$.  Thus $\RII$
Lem.~\texttt{lem:Kmax}, its interior maximiser $n_*$, the envelope $K^\star$ and
the gate of $\ZB$ play no role: $\tfrac{x^{m-3}(m-1-A)_+}{m}$ is handled
directly.
\end{remark}

\subsection*{Monotonicity toolkit}
Throughout $e\in(0,\tfrac13)$, $\kappa>0$, $D=1+e+2\kappa e$.

\begin{lemma}[quoted from $\ZB$]\label{zb9:mono}
\begin{enumerate}[label=\textup{(\roman*)},itemsep=1pt]
\item $x=(1-e)/D$ is strictly decreasing in $e$ and in $\kappa$.
\item $u=1-x$ is strictly increasing in $e$ and in $\kappa$.
\item $A=(1-e)(1+\kappa)/(\kappa D)$ is strictly decreasing in $e$ and in
$\kappa$.
\item $G(e):=\sqrt{1-e}\,(1-\lbar)(1-\lbar^{m-1})/(1+\lbar)$, with
$\lbar=\sqrt{2e/(1-e)}$, is positive and strictly decreasing on $(0,\tfrac13)$
for each fixed $m\ge2$ (product of positive decreasing factors).
\item $\gamma(e):=\tfrac7{25}-e$ is decreasing, $\kxi(e)$ increasing, and
$\kxi/(4\gamma)$ is strictly increasing on $[\tfrac1{60},\tfrac18]$ with maximum
$\kxi(\tfrac18)/(4\gamma(\tfrac18))=\tfrac{8/49}{31/50}=\tfrac{400}{1519}<1$
(the exact comparison $400<1519$ of $\ZB$).
\item $e\mapsto x(e,\gamma(e))$ is strictly decreasing on $[\tfrac1{60},\tfrac18]$;
$e\mapsto x(e,\kxi(e))$ is strictly decreasing on all of $[\tfrac1{60},e_1]$.
\item $\khi=\min(\kmx,\kq)$ is decreasing in $e$.
\end{enumerate}
\end{lemma}

\begin{lemma}[the $\gamma$-curve defect atom is monotone]\label{zb9:Agamma}
$e\mapsto A(e,\gamma(e))$ is strictly increasing on $[\tfrac1{60},\tfrac18]$.
\end{lemma}

\begin{proof}
On $\kappa=\gamma(e)=\tfrac7{25}-e$, $2\kappa e=\tfrac{14}{25}e-2e^2$, so
$D=1+\tfrac{39}{25}e-2e^2$ and, by \eqref{zb9:chart},
$A(e,\gamma(e))=\tfrac{(1-e)(\frac{32}{25}-e)}{(\frac7{25}-e)(1+\frac{39}{25}e-2e^2)}=:N/M$
with $N,M>0$ on $[\tfrac1{60},\tfrac18]$.  Then $A'$ has the sign of
$P:=N'M-NM'$, a fixed quartic with rational coefficients, $P>0$ on
$[\tfrac1{60},\tfrac18]$ (the script certifies $P>0$ at $6000$ rational points;
$A$ is monotone with endpoint values
$A(\tfrac1{60},\gamma)=\tfrac{3354150}{729091}=4.6004\ldots$,
$A(\tfrac18,\gamma)=\tfrac{3300}{589}=5.6027\ldots$).  Only
$A(e,\gamma(e))\ge A(\tfrac1{60},\gamma(\tfrac1{60}))=4.6004\ldots$ is used.
\end{proof}

\subsection*{Envelope: the $\Pi$, $\Lambda$, and defect brackets}

\begin{lemma}[the $\Pi$-envelope]\label{zb9:env}
On the strip $\kappa\ge\kxi(e)$:
\begin{enumerate}[label=\textup{(\alph*)},itemsep=1pt]
\item $\Pi\ge G(e)(1-\eps)$ with $\eps\le\min(\tfrac14,\tfrac{\kxi(e)}{4\kappa})$;
in particular $\Pi\ge\tfrac34G(e)$, and if $\kappa\ge\gamma(e)$ then
$\Pi\ge G(e)(1-\tfrac{\kxi(e)}{4\gamma(e)})$.
\item $G(e)\ge G_\downarrow(e_R)$ for $e\le e_R$, where, with $b:=e_R$ and
$5$-digit surrogates $c_\ell(b),c_s(b)$ satisfying (exact rational comparisons)
\begin{equation}\label{zb9:surr}
 c_\ell^2\ \ge\ \frac{2b}{1-b},\qquad c_s^2\ \le\ 1-b,
\end{equation}
one sets $G_\downarrow(b):=c_s\,\dfrac{(1-c_\ell)(1-(2b/(1-b))^{(m-1)/2})}{1+c_\ell}$.
\end{enumerate}
\end{lemma}

\begin{proof}
(a) The surrogate bounds $\ell\le\lbar$, $y\le\ell\le\lbar$,
$y^{m-1}\le\lbar^{m-1}$ (all from $L^2\le\alpha e$, $\RII$
Lem.~\texttt{elementary}) give $\Pi\ge G(e)(1-\eps)$.  Since $\rho\ge0$,
$\eps=\tfrac{e}{4(1-e)^2(1+\rho)\kappa}\le\tfrac{\kxi(e)}{4\kappa}$; and
$\kappa\ge\kxi$ gives $\eps\le\tfrac14$; if $\kappa\ge\gamma(e)$ then
$\tfrac{\kxi}{4\kappa}\le\tfrac{\kxi}{4\gamma}$.
(b) $(m-1)/2\in\{4,5,6\}$ is an integer, so
$\lbar^{m-1}=(2e/(1-e))^{(m-1)/2}\le(2b/(1-b))^{(m-1)/2}$; combine with
$\sqrt{1-e}\ge\sqrt{1-b}\ge c_s$, $\lbar\le c_\ell$, monotonicity of
$(1-t)/(1+t)$.
\end{proof}

\begin{lemma}[the $\Lambda$- and defect-brackets]\label{zb9:brk}
Fix $E=[e_L,e_R]$ and $c_\ell(e_R)$ as in \eqref{zb9:surr}.
\begin{enumerate}[label=\textup{(\alph*)},itemsep=1pt]
\item $\displaystyle
\Lambda\le\overline\Lambda(E):=\Bigl(\tfrac{1-e_L}{1+e_L}\Bigr)^{m-2}
\Bigl(\tfrac{2e_R}{1-e_R}\Bigr)^{(m-3)/2}c_\ell(e_R)\tfrac{(1-e_L)^2}{m\,e_L}$
for every strip point with $e\in E$ (using $\kappa\ge\kxi$).
\item For a box $E\times[\kappa_-,\kappa_+]$ on which $A\ge\underline A$ and
$x\le\overline x$ hold,
$\tfrac{x^{m-3}(m-1-A)_+}{m}\le D(\overline x,\underline A)
:=\tfrac{\overline x^{\,m-3}(m-1-\underline A)_+}{m}$.
\end{enumerate}
\end{lemma}

\begin{proof}
(a) $x\le\tfrac{1-e}{1+e}\le\tfrac{1-e_L}{1+e_L}$
(Lemma~\ref{zb9:mono}(i) with $\kappa\downarrow0$, then $e\ge e_L$);
$\tfrac1{m\kappa}\le\tfrac{(1-e)^2}{me}\le\tfrac{(1-e_L)^2}{me_L}$ (from
$\kappa\ge\kxi$); and
$(e/\alpha)^{m/2-1}=(2e/(1-e))^{(m-3)/2}\sqrt{2e/(1-e)}\le(2e_R/(1-e_R))^{(m-3)/2}c_\ell(e_R)$.
Multiply.  (b) $x^{m-3}$ increasing in $x$, $(m-1-A)_+$ decreasing in $A$.
\end{proof}

\subsection*{The battle: three cases and the displayed tables}
Fix $m\in\{9,11,13\}$.  For each row $E=[e_L,e_R]$ the $m$-table lists a lower
bound $L(E)$ for $\Pi$ and an upper bound $R(E)$ for
$\tfrac{x^{m-3}(m-1-A)_+}{m}+\Lambda$; the row asserts $L(E)\ge R(E)$.  We
abbreviate $D(\overline x,\underline A)=\overline x^{\,m-3}(m-1-\underline A)_+/m$.

\begin{proposition}[Case S: $\kappa\le\gamma(e)$, $e\le\tfrac18$]\label{zb9:caseS}
Let $e\in[\tfrac1{60},\tfrac18]$, $\kxi(e)\le\kappa\le\gamma(e)$.  Then
\eqref{zb9:reduced} holds, since for the Case-S block $E=[e_L,e_R]$ containing
$e$,
\[
 \Pi\ \ge\ \underbrace{\tfrac34 G_\downarrow(e_R)}_{=:L_S(E)},\qquad
 \frac{x^{m-3}(m-1-A)_+}{m}+\Lambda\ \le\
 \underbrace{D\bigl(x(e_L,\kxi(e_L)),A(e_L,\gamma(e_L))\bigr)+\overline\Lambda(E)}_{=:R_S(E)},
\]
and every block certifies $L_S(E)\ge R_S(E)$.
\end{proposition}

\begin{proof}
$\Pi\ge\tfrac34G(e)\ge\tfrac34G_\downarrow(e_R)$ by
Lemma~\ref{zb9:env}(a),(b).  For the defect box:
$x\le x(e,\kxi(e))\le x(e_L,\kxi(e_L))$ (Lemma~\ref{zb9:mono}(i),(vi)) and
$A\ge A(e,\gamma(e))\ge A(e_L,\gamma(e_L))$ (Lemma~\ref{zb9:mono}(iii) as
$\kappa\le\gamma(e)$, then Lemma~\ref{zb9:Agamma}), so
$\underline A=A(e_L,\gamma(e_L))$, $\overline x=x(e_L,\kxi(e_L))$ are valid;
$\Lambda\le\overline\Lambda(E)$ is Lemma~\ref{zb9:brk}(a).
\end{proof}

\begin{proposition}[Case L$\gamma$: $\kappa\ge\gamma(e)$, $e\le\tfrac18$]\label{zb9:caseLg}
Let $e\in[\tfrac1{60},\tfrac18]$, $\gamma(e)\le\kappa\le\khi(e)$.  Then
\eqref{zb9:reduced} holds, since for the block $E$ containing $e$,
\[
 \Pi\ \ge\ \underbrace{G_\downarrow(e_R)\bigl(1-\tfrac{\kxi(e_R)}{4\gamma(e_R)}\bigr)}_{=:L_{L\gamma}(E)},\qquad
 \frac{x^{m-3}(m-1-A)_+}{m}+\Lambda\ \le\
 \underbrace{D\bigl(x(e_L,\gamma(e_L)),A(e_R,\khi(e_L))\bigr)+\overline\Lambda(E)}_{=:R_{L\gamma}(E)},
\]
and every block certifies $L_{L\gamma}(E)\ge R_{L\gamma}(E)$.
\end{proposition}

\begin{proof}
$\Pi\ge G(e)(1-\kxi(e)/(4\gamma(e)))$ by Lemma~\ref{zb9:env}(a)
($\kappa\ge\gamma(e)$); the map is a product of two positive decreasing factors
(Lemma~\ref{zb9:mono}(iv),(v)), hence $\ge L_{L\gamma}(E)$ at $e=e_R$, and
$G(e_R)\ge G_\downarrow(e_R)$.  Box: $x\le x(e,\gamma(e))\le x(e_L,\gamma(e_L))$;
and $A\ge A(e_R,\khi(e_L))$ since $A$ decreases in $e,\kappa$ while
$\kappa\le\khi(e)\le\khi(e_L)$ and $e\le e_R$.  Then Lemma~\ref{zb9:brk}.
\end{proof}

\begin{proposition}[Case L$\xi$: $e\ge\tfrac18$]\label{zb9:caseLx}
Let $e\in[\tfrac18,e_1]$, $\kxi(e)\le\kappa\le\khi(e)$.  Then
\eqref{zb9:reduced} holds, since for the block $E$ containing $e$,
\[
 \Pi\ \ge\ \underbrace{\tfrac34 G_\downarrow(e_R)}_{=:L_{L\xi}(E)},\qquad
 \frac{x^{m-3}(m-1-A)_+}{m}+\Lambda\ \le\
 \underbrace{D\bigl(x(e_L,\kxi(e_L)),A(e_R,\khi(e_L))\bigr)+\overline\Lambda(E)}_{=:R_{L\xi}(E)},
\]
and every block certifies $L_{L\xi}(E)\ge R_{L\xi}(E)$.
\end{proposition}

\begin{proof}
Left side as in Case~S ($\Pi\ge\tfrac34G_\downarrow(e_R)$); defect box as in
Case~L$\gamma$ ($\overline x=x(e_L,\kxi(e_L))$, $\underline A=A(e_R,\khi(e_L))$);
$\Lambda\le\overline\Lambda(E)$.
\end{proof}

\begin{proposition}[the fixed-$m$ battle]\label{zb9:battle}
For each $m\in\{9,11,13\}$, \eqref{zb9:reduced}---hence \eqref{eq:target}---holds
at every strip point $e\in[\tfrac1{60},e_1]$, $\kappa\in[\kxi(e),\khi(e)]$.
\end{proposition}

\begin{proof}
If $R_m\le0$, \eqref{eq:target} is trivial; assume $R_m>0$.  For $e\le\tfrac18$
every $\kappa$ is Case~S ($\kappa\le\gamma(e)$) or Case~L$\gamma$
($\kappa\ge\gamma(e)$); for $e\ge\tfrac18$ it is Case~L$\xi$.  Each case is
closed by the corresponding block of Table~\ref{zb9:tab9} ($m=9$),
\ref{zb9:tab11} ($m=11$), or \ref{zb9:tab13} ($m=13$), whose blocks tile
$[\tfrac1{60},\tfrac18]$ (S and L$\gamma$) and $[\tfrac18,e_1]$ (L$\xi$); every
row satisfies $L\ge R$.  Then Lemma~\ref{zb9:red} gives \eqref{eq:target}.
\end{proof}

\begin{table}[h]\centering
\caption{$m=9$.  Row certificate $L(E)\ge R(E)$; $L$ rounded down, $R$ rounded
up.  Blocks S and L$\gamma$ tile $[\tfrac1{60},\tfrac18]$; block L$\xi$ tiles
$[\tfrac18,e_1]$.}\label{zb9:tab9}
\begin{tabular}{clcc}\toprule
case & $E=[e_L,e_R]$ & $L(E)\ge$ & $R(E)\le$\\\midrule
S       & $[1/60,\,7/100]$   & $0.3187$ & $0.3150$\\
S       & $[7/100,\,1/8]$    & $0.2113$ & $0.1470$\\\midrule
L$\gamma$ & $[1/60,\,7/200]$ & $0.5437$ & $0.5380$\\
L$\gamma$ & $[7/200,\,3/50]$ & $0.4235$ & $0.4206$\\
L$\gamma$ & $[3/50,\,1/11]$  & $0.3107$ & $0.3035$\\
L$\gamma$ & $[1/11,\,3/25]$  & $0.2218$ & $0.2049$\\
L$\gamma$ & $[3/25,\,1/8]$   & $0.2076$ & $0.1404$\\\midrule
L$\xi$  & $[1/8,\,17/100]$   & $0.1457$ & $0.1378$\\
L$\xi$  & $[17/100,\,1/5]$   & $0.1078$ & $0.0619$\\
L$\xi$  & $[1/5,\,2033/10000]$ & $0.1039$ & $0.0323$\\
\bottomrule\end{tabular}
\end{table}

\begin{table}[h]\centering
\caption{$m=11$.  Layout as in Table~\ref{zb9:tab9}.}\label{zb9:tab11}
\begin{tabular}{clcc}\toprule
case & $E=[e_L,e_R]$ & $L(E)\ge$ & $R(E)\le$\\\midrule
S       & $[1/60,\,1/21]$    & $0.3802$ & $0.3744$\\
S       & $[1/21,\,3/25]$    & $0.2204$ & $0.2177$\\
S       & $[3/25,\,1/8]$     & $0.2123$ & $0.0460$\\\midrule
L$\gamma$ & $[1/60,\,1/28]$  & $0.5397$ & $0.5347$\\
L$\gamma$ & $[1/28,\,7/100]$ & $0.3842$ & $0.3772$\\
L$\gamma$ & $[7/100,\,1/8]$  & $0.2086$ & $0.2061$\\\midrule
L$\xi$  & $[1/8,\,1/5]$      & $0.1114$ & $0.0831$\\
L$\xi$  & $[1/5,\,2033/10000]$ & $0.1076$ & $0.0126$\\
\bottomrule\end{tabular}
\end{table}

\begin{table}[h]\centering
\caption{$m=13$.  Layout as in Table~\ref{zb9:tab9}.}\label{zb9:tab13}
\begin{tabular}{clcc}\toprule
case & $E=[e_L,e_R]$ & $L(E)\ge$ & $R(E)\le$\\\midrule
S       & $[1/60,\,1/26]$    & $0.4111$ & $0.4056$\\
S       & $[1/26,\,1/10]$    & $0.2555$ & $0.2533$\\
S       & $[1/10,\,1/8]$     & $0.2126$ & $0.0565$\\\midrule
L$\gamma$ & $[1/60,\,1/26]$  & $0.5245$ & $0.5181$\\
L$\gamma$ & $[1/26,\,17/200]$& $0.3307$ & $0.3125$\\
L$\gamma$ & $[17/200,\,1/8]$ & $0.2089$ & $0.1096$\\\midrule
L$\xi$  & $[1/8,\,1/5]$      & $0.1132$ & $0.0478$\\
L$\xi$  & $[1/5,\,2033/10000]$ & $0.1095$ & $0.0049$\\
\bottomrule\end{tabular}
\end{table}

\section{Zone C: the $\xi\le1$ slice, including the Tur\'an corner}\label{sec:zoneC}

This section proves \eqref{eq:target} on the whole $\xi\le1$ slice
(equivalently $d\le e^2/(1-e)^2$) by specialising the verified analytic Zone-C
proof $\ZC$ to the fixed exponent $n:=m-2\in\{7,9,11\}$.  Write
$\Def_m:=R_m/(\alpha^3p^{\,n})$.  The single new ingredient beyond $\ZC$ is a
\emph{mean-value Case-II row bound} (\S below) that removes the $m$-uniform
supremum loss of $\ZC$ for the small exponents; with it every range closes on a
\emph{fixed, fully displayed} interval cover, with no adaptive subdivision.

\begin{theorem}\label{zc9:thm:main}
Let $m\in\{9,11,13\}$ and let $(q,\alpha)$ be any admissible point with $\xi\le1$
and $q>\tfrac13$.  Then \eqref{eq:target} holds.
\end{theorem}

The proof partitions $e\in(0,\tfrac13)$ into the pieces of
Table~\ref{zc9:tab:split}; $e^\star_m=\tfrac13-2\Delta_m$ with
$\Delta_{13}=\tfrac{13}{600},\Delta_{11}=\tfrac9{600},\Delta_9=\tfrac7{600}$.
The strip end $\bar e_m$ equals the corner start $e^\star_m$ for $m\in\{11,13\}$;
for $m=9$ a short sliver bridges them.  Here
$\bar e_{13}=e^\star_{13}=\tfrac{29}{100}$,
$\bar e_{11}=e^\star_{11}=\tfrac{91}{300}$, and $\bar e_{9}=\tfrac{29}{100}$
with $e^\star_9=\tfrac{31}{100}$.
\begin{table}[h]\centering\small
\caption{Zone-C partition.}\label{zc9:tab:split}
\begin{tabular}{llll}\toprule
piece & range & method & where\\\midrule
Zone A & $0<e\le\tfrac1{60}$ & $R_m\le0$ & Cor.~\ref{zc9:cor:za}\\
strip  & $\tfrac1{60}\le e\le \bar e_m$ & fixed $\varphi$-row table & Prop.~\ref{zc9:prop:strip}\\
sliver & ($m=9$) $\tfrac{29}{100}\le e\le\tfrac{31}{100}$ & fixed $(e,t)$ cell table & Prop.~\ref{zc9:prop:sliver}\\
corner & $e^\star_m\le e<\tfrac13$ & dual-certificate chain & Prop.~\ref{zc9:prop:wide}\\
\bottomrule\end{tabular}
\end{table}

\subsection*{What is inherited verbatim from $\ZC$}
$\ZC$ uses $m\ge15$ \emph{only} through $n\ge13$ (the exponent in $G_2$),
$1-y^{m-1}\ge1-y^{14}$, and the constant $15$ in its Lemma~\texttt{master};
replacing $13\rightsquigarrow n=m-2$, $14\rightsquigarrow m-1$,
$15\rightsquigarrow m$ leaves the following intact for every fixed $m\ge9$
($n\ge7$).  With
$\lambda_x=\ln\tfrac p\alpha$, $\Delta=\ln\tfrac\alpha q$,
$\lambda_s=\lambda_x+\Delta$, $G_2=\ell^2(1-(y/s)^{n})$, $G=\ln(1+G_2)$,
$n_g=G/\Delta$:
\begin{itemize}
\item \textbf{$\ZC$ Lem.~\texttt{majorant}}:
 $\alpha\Def_m=x^n-s^n-\ell^2(s^n-y^n)\le\varphi(n):=x^n-(1+G_2)s^n$.
\item \textbf{$\ZC$ Lem.~\texttt{sup}}: $\varphi(t)=x^t-(1+G_2)s^t$ is unimodal
 with $\sup_t\varphi=(\Delta/\lambda_s)h\eu^{-w}$, $w:=\lambda_xG/\Delta$,
 $h:=\exp(-\tfrac{\lambda_x}\Delta\ln(1+\tfrac\Delta{\lambda_x}))\in(\eu^{-1},1)$,
 $h\le\eu^{-1+v/2}$, $v=\Delta/\lambda_x$.
\item \textbf{$\ZC$ Lem.~\texttt{scalar},~\texttt{payfloor},~\texttt{master}}:
 with $c_e=\tfrac{2(1-x^{m-1})(1-y^{m-1})}{(1+x)(1+y)}$,
 $K_2=\tfrac{\sqrt{2\alpha}(1-y^{m-1})}{2\alpha^3(1+y)}$, the payment floors are
 $C_m\psi/(\alpha^3p^{n})\ge m c_e\kappa^2$ if $2\rho\xi\le1$ and $\ge mK_2fd$ if
 $2\rho\xi>1$; and, if $R_m>0$, \eqref{eq:target} follows from the applicable
 \emph{Case-I} master inequality
 \begin{align}
  &\textup{(deep,}\ 2\rho\xi\le1,\ n_g\ge m)&
   h\eu^{-w}e^2&\le\lambda_x\alpha c_e q^2G,\label{zc9:eq:SCdeep}\\
  &\textup{(shallow,}\ 2\rho\xi\le1,\ n_g<m)&
   h\eu^{-w}e^2&\le m\lambda_x\alpha c_e q^2\Delta.\label{zc9:eq:SCsh}
 \end{align}
\item \textbf{$\ZC$ Lem.~\texttt{bracket}}: over $E=[e_L,e_R]$ with
 $d\le d_{\max}(e)=\min(e^2/(1-e)^2,\delta)$, the monotone-endpoint constants
 $\bar q,D,\Lambda,\lambda_L,\widehat G_2,Y,\gamma,a,w^-,v^+,h^+,x_U,y_U,c^-,f^-,K^-$
 (recorded in \S\ref{zc9:sec:strip}) satisfy $q\ge\bar q$, $\Delta\le\Lambda$,
 $\lambda_x\ge\lambda_L$, $\lambda_xG\ge a$, $w\ge w^-$, $h\le h^+$,
 $c_e\ge c^-$, $f\ge f^-$, $K_2\ge K^-$, $x\le x_U$, $y\le y_U$.
\end{itemize}
The oddness of $m$ and $\alpha\le r(q)$ enter only through $L^2>0$,
$0\le\ell<1$, $0<y<s$, $L<q$ ($\RII$ Lem.~\texttt{frontier-gap}).  We replace
only the \emph{Case-II} treatment of $\ZC$; its Case-II master inequality,
retained as one of two admissible sufficient conditions, is
\begin{equation}\label{zc9:eq:SCII-old}
 \textup{(Case II, }2\rho\xi>1):\qquad h\eu^{-w}\ \le\ m\lambda_x\alpha K_2 f q .
\end{equation}

\begin{corollary}[Zone A restated]\label{zc9:cor:za}
For every $m\in\{9,11,13\}$ and every admissible point with $\xi\le1$ and
$e\le\tfrac1{60}$ one has $R_m\le0$; hence \eqref{eq:target} holds trivially.
\end{corollary}

\begin{proof}
By Corollary~\ref{cor:gate-smalle}, $R_m>0$ and $e\le\tfrac1{60}$ force
$\xi\ge3.68>1$; so on $\{e\le\tfrac1{60},\ \xi\le1\}$ necessarily $R_m\le0$.
(Part~7 of the validation script confirms $\max R_m<0$ there.)
\end{proof}

\subsection*{A mean-value Case-II row bound}
The only structural change from $\ZC$ is to replace, in Case~II, the majorant
$\alpha\Def_m\le\varphi(n)\le\sup_t\varphi$ (which loses a factor $\approx4$ at
$n=7$) by a sharper mean-value majorant.

\begin{lemma}[mean-value defect]\label{zc9:lem:mvt}
On all of $\D$, $\alpha\,\Def_m\le\dfrac{n\,x^{\,n-1}}p\,N$ where
$N:=d-\ell^{\,m-1}(q-L)$, and $0<N\le d$ whenever $R_m>0$.
\end{lemma}

\begin{proof}
By $\ZC$ Lem.~\texttt{majorant}, $\alpha\Def_m=(x^n-s^n)-\ell^2(s^n-y^n)$.  The
mean value theorem gives $c'\in(s,x)$, $c\in(y,s)$ with
$x^n-s^n=nc'^{\,n-1}(x-s)\le nx^{n-1}(x-s)$ and
$s^n-y^n=nc^{\,n-1}(s-y)\ge ny^{n-1}(s-y)$.  Since $\ell^2>0$,
$\alpha\Def_m\le nx^{n-1}(x-s)-\ell^2 ny^{n-1}(s-y)$.  Now $x-s=d/p$,
$s-y=(q-L)/p$, and $y=\ell x$ gives
$\ell^2 y^{n-1}=\ell^{\,n+1}x^{n-1}=\ell^{\,m-1}x^{n-1}$ (as $m-1=n+1$).  Hence
$\alpha\Def_m\le\tfrac{n x^{n-1}}p(d-\ell^{\,m-1}(q-L))=\tfrac{n x^{n-1}}p N$.
Finally $q>L$ on $\D$ ($L<q\iff q>\tfrac13$) and $\ell>0$, so $N\le d$; and
$R_m>0$ forces $\alpha\Def_m>0$, whence $N>0$.  (Part~4 of the script verifies
both on a dense grid.)
\end{proof}

\begin{lemma}[Case-II master, MVT form]\label{zc9:lem:mvtII}
Let $R_m>0$ and $2\rho\xi>1$.  Then \eqref{eq:target} holds as soon as
\begin{equation}\label{zc9:eq:SCII-mvt}
 \frac{n\,x^{\,n-1}}{p\,\alpha}\ \le\ m\,K_2\,f .
\end{equation}
\end{lemma}

\begin{proof}
By $\ZC$ Lem.~\texttt{payfloor}, $2\rho\xi>1$ gives
$C_m\psi/(\alpha^3p^{n})\ge mK_2fd$, so it suffices that $\Def_m\le mK_2fd$.  By
Lemma~\ref{zc9:lem:mvt}, $\Def_m\le\tfrac{n x^{n-1}}{p\alpha}N\le
\tfrac{n x^{n-1}}{p\alpha}d$; dividing by $d>0$ gives \eqref{zc9:eq:SCII-mvt}.
\end{proof}

\begin{lemma}[assembled Case-II row bound]\label{zc9:lem:mvtrow}
Over $E=[e_L,e_R]$ with $d\le d_{\max}$, put $p^-:=\tfrac{1+e_L}2$ and take
$x_U,f^-,K^-$ from $\ZC$ Lem.~\texttt{bracket}.  If
$\Psi_{\mathrm{II}}^{\mathrm{mvt}}(E):=\tfrac{n\,x_U^{\,n-1}}{p^-\,\alpha(e_R)\,m\,K^-\,f^-}\le1$,
then \eqref{zc9:eq:SCII-mvt} holds throughout $E\cap\{2\rho\xi>1,R_m>0\}$.
\end{lemma}

\begin{proof}
$x\le x_U$ gives $nx^{n-1}\le nx_U^{\,n-1}$; also $p\ge\tfrac{1+e}2\ge p^-$,
$\alpha\ge\alpha(e_R)$, $f\ge f^-$, $K_2\ge K^-$ ($\ZC$ Lem.~\texttt{bracket}).
Hence $\tfrac{n x^{n-1}}{p\alpha}\le\tfrac{n x_U^{\,n-1}}{p^-\alpha(e_R)}$ and
$mK_2f\ge mK^-f^-$.
\end{proof}

Since \eqref{zc9:eq:SCII-old} and \eqref{zc9:eq:SCII-mvt} are both sufficient
for Case~II, the effective Case-II row bound on $E$ is
\begin{equation}\label{zc9:eq:PsiII}
 \Psi_{\mathrm{II}}(E):=\min\bigl(\Psi_{\mathrm{II}}^{\mathrm{[ZC]}}(E),\
 \Psi_{\mathrm{II}}^{\mathrm{mvt}}(E)\bigr),\qquad
 \Psi_{\mathrm{II}}^{\mathrm{[ZC]}}(E)=\frac{h^+\eu^{-a/\Lambda}}
 {m\lambda_L\alpha(e_R)K^-f^-\bar q}.
\end{equation}
The $\ZC$ form is tighter at small $e$ (where $N\ll d$); the MVT form is tighter
on the ridge.  No $d$-lower-bound is needed.

\subsection{The strip $\tfrac1{60}\le e\le\bar e_m$}\label{zc9:sec:strip}

Fix $E=[e_L,e_R]$ with $d\le d_{\max}(e)$.  The $\ZC$ Lem.~\texttt{bracket}
constants (elementary evaluations at the rational endpoints;
$\alpha(u)=\tfrac{1-u}2$, $\delta(u)=\tfrac{1-3u}6$) are
\[
 \bar q=\max\bigl(\tfrac13,\alpha(e_R)-\tfrac{e_R^2}{(1-e_R)^2}\bigr),\ \
 D=\min\bigl(\tfrac{e_R^2}{(1-e_R)^2},\delta(e_L)\bigr),\ \
 \Lambda=\tfrac D{\bar q},\ \
 \lambda_L=\ln\tfrac{1+e_L}{1-e_L},\ \
 \widehat G_2=\tfrac{2e_R}{1-e_R},
\]
\[
 Y=\tfrac{\sqrt{e_R(1-e_R)/2}}{\bar q},\quad
 \gamma=\Bigl[\tfrac2{1-e_L}-\tfrac{D(e_R+D)}{\alpha(e_R)^2e_L}\Bigr]
 (1-Y^n)\tfrac{\ln(1+\widehat G_2)}{\widehat G_2},\quad a=\lambda_L\gamma e_L,
\]
\[
 w^-=2\gamma\max\bigl((1-e_R)^2\bar q,\tfrac{e_L^2}{3\delta(e_L)}\bigr),\quad
 v^+=\min\bigl(\tfrac{e_R}{2(1-e_R)^2\bar q},\tfrac{3\delta(e_L)}{\lambda_L}\bigr),\quad
 h^+=\eu^{-1+v^+/2},
\]
\[
 x_U=\tfrac{1-e_L}{1+e_L},\ \
 y_U=\tfrac{\sqrt{e_R(1-e_R)/2}}{\alpha(e_R)+e_R},\ \
 c^-=\tfrac{2(1-x_U^{m-1})(1-y_U^{m-1})}{(1+x_U)(1+y_U)},\ \
 f^-=\alpha(e_R)-\sqrt{\tfrac{e_R(1-e_R)}2},
\]
\[
 K^-=\tfrac{\sqrt{2\alpha(e_R)}(1-y_U^{m-1})}{2\alpha(e_L)^3(1+y_U)},\qquad
 p^-=\tfrac{1+e_L}2 .
\]
As in $\ZC$ Lem.~\texttt{rows}, the Case-I conditions
\eqref{zc9:eq:SCdeep}--\eqref{zc9:eq:SCsh} reduce, via Lem.~\texttt{bracket}, to
$\Psi_{\mathrm{deep}}(E)\le1$ and (when $\Lambda>\gamma e_L/m$)
$\Psi_{\mathrm{sh}}(E)\le1$, where
\[
 \Psi_{\mathrm{deep}}=\frac{h^+\eu^{-\max(w^-,m\lambda_L)}}
 {2\gamma\,\alpha(e_R)\,c^-\,\bar q^{\,2}},\qquad
 \Psi_{\mathrm{sh}}=\frac{h^+(e_R/2)\eu^{-a/\Delta^*}}
 {m\,\alpha(e_R)\,c^-\,\bar q^{\,2}\,\Delta^*},\quad
 \Delta^*=\min\bigl(\Lambda,\max(a,\tfrac{\gamma e_L}m)\bigr),
\]
and Case~II reduces to $\Psi_{\mathrm{II}}(E)\le1$ of \eqref{zc9:eq:PsiII}.
Thus if $\max(\Psi_{\mathrm{deep}},\Psi_{\mathrm{sh}},\Psi_{\mathrm{II}})\le1$ on
$E$ then \eqref{eq:target} holds at every point of $\D$ with $e\in E$ and
$d\le d_{\max}(e)$.

\begin{proposition}[strip]\label{zc9:prop:strip}
For each $m\in\{9,11,13\}$, \eqref{eq:target} holds at every admissible point
with $\xi\le1$ and $\tfrac1{60}\le e\le\bar e_m$, where
$\bar e_{13}=\bar e_9=\tfrac{29}{100}$ and $\bar e_{11}=\tfrac{91}{300}$.
\end{proposition}

\begin{proof}
The cover of $[\tfrac1{60},\bar e_m]$ is the \textbf{fixed} rational partition of
Tables~\ref{zc9:tab:m13}--\ref{zc9:tab:m9}: the $\ZC$ breakpoints
$\tfrac1{60},\tfrac1{48},\dots,\tfrac{29}{100}$ (extended by
$\tfrac3{10},\tfrac{91}{300}$ for $m=11$), with three bisected for $m=9$ and one
for $m=11$ (midpoints $\tfrac3{160},\tfrac{27}{100},\tfrac{57}{200}$ displayed as
ordinary rows).  Every row lists an upper bound (rounded up) for the
corresponding $\Psi$; all are $<1$.  On the first rows $\Lambda\le\gamma e_L/m$,
so the shallow case is empty and $\Psi_{\mathrm{sh}}$ is marked ``$-$''.
Part~3 recomputes every row; Part~2 independently verifies, on a dense
$(e,\kappa)$ grid over each interval, that the tabulated $\Psi$ of the
\emph{applicable} branch dominates the true ratio $\Def_m/\text{floor}$.
\end{proof}
\begin{table}[p]\centering\small
\caption{$m=13$ strip: $20$ rows on $[\tfrac1{60},\tfrac{29}{100}]$ (the [ZC] cover, unbisected).}\label{zc9:tab:m13}
\begin{tabular}{llcccl}\toprule
$e_L$&$e_R$&$\Psi_{\mathrm{deep}}\!\le$&$\Psi_{\mathrm{sh}}\!\le$&$\Psi_{\mathrm{II}}\!\le$&bind\\\midrule
$\tfrac{1}{60}$&$\tfrac{1}{48}$&$0.442$&$-$&$0.805$&II\\
$\tfrac{1}{48}$&$\tfrac{1}{40}$&$0.390$&$-$&$0.619$&II\\
$\tfrac{1}{40}$&$\tfrac{1}{32}$&$0.370$&$-$&$0.628$&II\\
$\tfrac{1}{32}$&$\tfrac{1}{25}$&$0.355$&$-$&$0.601$&II\\
$\tfrac{1}{25}$&$\tfrac{1}{20}$&$0.344$&$-$&$0.508$&II\\
$\tfrac{1}{20}$&$\tfrac{1}{16}$&$0.356$&$0.618$&$0.479$&s\\
$\tfrac{1}{16}$&$\tfrac{1}{13}$&$0.337$&$0.477$&$0.447$&s\\
$\tfrac{1}{13}$&$\tfrac{1}{10}$&$0.252$&$0.445$&$0.550$&II\\
$\tfrac{1}{10}$&$\tfrac{1}{8}$&$0.153$&$0.367$&$0.549$&II\\
$\tfrac{1}{8}$&$\tfrac{3}{20}$&$0.095$&$0.335$&$0.435$&II\\
$\tfrac{3}{20}$&$\tfrac{17}{100}$&$0.059$&$0.316$&$0.284$&s\\
$\tfrac{17}{100}$&$\tfrac{9}{50}$&$0.038$&$0.285$&$0.190$&s\\
$\tfrac{9}{50}$&$\tfrac{19}{100}$&$0.034$&$0.315$&$0.164$&s\\
$\tfrac{19}{100}$&$\tfrac{1}{5}$&$0.032$&$0.368$&$0.142$&s\\
$\tfrac{1}{5}$&$\tfrac{21}{100}$&$0.028$&$0.399$&$0.124$&s\\
$\tfrac{21}{100}$&$\tfrac{11}{50}$&$0.022$&$0.378$&$0.108$&s\\
$\tfrac{11}{50}$&$\tfrac{6}{25}$&$0.019$&$0.427$&$0.109$&s\\
$\tfrac{6}{25}$&$\tfrac{13}{50}$&$0.012$&$0.429$&$0.088$&s\\
$\tfrac{13}{50}$&$\tfrac{7}{25}$&$0.009$&$0.481$&$0.077$&s\\
$\tfrac{7}{25}$&$\tfrac{29}{100}$&$0.006$&$0.454$&$0.058$&s\\
\bottomrule\end{tabular}\end{table}
\begin{table}[p]\centering\small
\caption{$m=11$ strip: $23$ rows on $[\tfrac1{60},\tfrac{91}{300}]$ ([ZC] cover $+$ tail $\tfrac3{10}\!\to\!\tfrac{91}{300}$; $[\tfrac1{60},\tfrac1{48}]$ bisected at $\tfrac3{160}$).}\label{zc9:tab:m11}
\begin{tabular}{llcccl}\toprule
$e_L$&$e_R$&$\Psi_{\mathrm{deep}}\!\le$&$\Psi_{\mathrm{sh}}\!\le$&$\Psi_{\mathrm{II}}\!\le$&bind\\\midrule
$\tfrac{1}{60}$&$\tfrac{3}{160}$&$0.496$&$-$&$0.688$&II\\
$\tfrac{3}{160}$&$\tfrac{1}{48}$&$0.462$&$-$&$0.609$&II\\
$\tfrac{1}{48}$&$\tfrac{1}{40}$&$0.450$&$-$&$0.732$&II\\
$\tfrac{1}{40}$&$\tfrac{1}{32}$&$0.425$&$-$&$0.742$&II\\
$\tfrac{1}{32}$&$\tfrac{1}{25}$&$0.403$&$-$&$0.710$&II\\
$\tfrac{1}{25}$&$\tfrac{1}{20}$&$0.385$&$-$&$0.601$&II\\
$\tfrac{1}{20}$&$\tfrac{1}{16}$&$0.393$&$0.807$&$0.566$&s\\
$\tfrac{1}{16}$&$\tfrac{1}{13}$&$0.418$&$0.614$&$0.528$&s\\
$\tfrac{1}{13}$&$\tfrac{1}{10}$&$0.368$&$0.565$&$0.650$&II\\
$\tfrac{1}{10}$&$\tfrac{1}{8}$&$0.242$&$0.457$&$0.650$&II\\
$\tfrac{1}{8}$&$\tfrac{3}{20}$&$0.164$&$0.414$&$0.695$&II\\
$\tfrac{3}{20}$&$\tfrac{17}{100}$&$0.113$&$0.391$&$0.502$&II\\
$\tfrac{17}{100}$&$\tfrac{9}{50}$&$0.079$&$0.354$&$0.365$&II\\
$\tfrac{9}{50}$&$\tfrac{19}{100}$&$0.075$&$0.397$&$0.328$&s\\
$\tfrac{19}{100}$&$\tfrac{1}{5}$&$0.076$&$0.473$&$0.296$&s\\
$\tfrac{1}{5}$&$\tfrac{21}{100}$&$0.070$&$0.520$&$0.268$&s\\
$\tfrac{21}{100}$&$\tfrac{11}{50}$&$0.056$&$0.497$&$0.244$&s\\
$\tfrac{11}{50}$&$\tfrac{6}{25}$&$0.052$&$0.571$&$0.256$&s\\
$\tfrac{6}{25}$&$\tfrac{13}{50}$&$0.037$&$0.585$&$0.227$&s\\
$\tfrac{13}{50}$&$\tfrac{7}{25}$&$0.029$&$0.666$&$0.215$&s\\
$\tfrac{7}{25}$&$\tfrac{29}{100}$&$0.020$&$0.654$&$0.176$&s\\
$\tfrac{29}{100}$&$\tfrac{3}{10}$&$0.020$&$0.787$&$0.189$&s\\
$\tfrac{3}{10}$&$\tfrac{91}{300}$&$0.016$&$0.728$&$0.170$&s\\
\bottomrule\end{tabular}\end{table}
\begin{table}[p]\centering\small
\caption{$m=9$ strip: $23$ rows on $[\tfrac1{60},\tfrac{29}{100}]$ ([ZC] cover; $[\tfrac1{60},\tfrac1{48}],[\tfrac{13}{50},\tfrac7{25}],[\tfrac7{25},\tfrac{29}{100}]$ bisected at $\tfrac3{160},\tfrac{27}{100},\tfrac{57}{200}$).}\label{zc9:tab:m9}
\begin{tabular}{llcccl}\toprule
$e_L$&$e_R$&$\Psi_{\mathrm{deep}}\!\le$&$\Psi_{\mathrm{sh}}\!\le$&$\Psi_{\mathrm{II}}\!\le$&bind\\\midrule
$\tfrac{1}{60}$&$\tfrac{3}{160}$&$0.601$&$-$&$0.841$&II\\
$\tfrac{3}{160}$&$\tfrac{1}{48}$&$0.557$&$-$&$0.745$&II\\
$\tfrac{1}{48}$&$\tfrac{1}{40}$&$0.541$&$-$&$0.894$&II\\
$\tfrac{1}{40}$&$\tfrac{1}{32}$&$0.507$&$-$&$0.907$&II\\
$\tfrac{1}{32}$&$\tfrac{1}{25}$&$0.477$&$-$&$0.867$&II\\
$\tfrac{1}{25}$&$\tfrac{1}{20}$&$0.449$&$-$&$0.734$&II\\
$\tfrac{1}{20}$&$\tfrac{1}{16}$&$0.452$&$-$&$0.693$&II\\
$\tfrac{1}{16}$&$\tfrac{1}{13}$&$0.474$&$0.849$&$0.646$&s\\
$\tfrac{1}{13}$&$\tfrac{1}{10}$&$0.558$&$0.770$&$0.798$&II\\
$\tfrac{1}{10}$&$\tfrac{1}{8}$&$0.396$&$0.613$&$0.803$&II\\
$\tfrac{1}{8}$&$\tfrac{3}{20}$&$0.295$&$0.552$&$0.868$&II\\
$\tfrac{3}{20}$&$\tfrac{17}{100}$&$0.226$&$0.525$&$0.874$&II\\
$\tfrac{17}{100}$&$\tfrac{9}{50}$&$0.173$&$0.479$&$0.691$&II\\
$\tfrac{9}{50}$&$\tfrac{19}{100}$&$0.174$&$0.546$&$0.647$&II\\
$\tfrac{19}{100}$&$\tfrac{1}{5}$&$0.187$&$0.667$&$0.608$&s\\
$\tfrac{1}{5}$&$\tfrac{21}{100}$&$0.183$&$0.744$&$0.574$&s\\
$\tfrac{21}{100}$&$\tfrac{11}{50}$&$0.154$&$0.716$&$0.545$&s\\
$\tfrac{11}{50}$&$\tfrac{6}{25}$&$0.151$&$0.840$&$0.597$&s\\
$\tfrac{6}{25}$&$\tfrac{13}{50}$&$0.118$&$0.874$&$0.574$&s\\
$\tfrac{13}{50}$&$\tfrac{27}{100}$&$0.083$&$0.794$&$0.484$&s\\
$\tfrac{27}{100}$&$\tfrac{7}{25}$&$0.079$&$0.874$&$0.497$&s\\
$\tfrac{7}{25}$&$\tfrac{57}{200}$&$0.068$&$0.869$&$0.467$&s\\
$\tfrac{57}{200}$&$\tfrac{29}{100}$&$0.068$&$0.938$&$0.484$&s\\
\bottomrule\end{tabular}\end{table}

\subsection{The $m=9$ junction sliver}\label{zc9:sec:sliver}

For $m=9$ the strip stops at $\bar e_9=\tfrac{29}{100}$ (beyond which the
$m$-uniform shallow bound $\Psi_{\mathrm{sh}}$ exceeds $1$) and the corner
starts at $e^\star_9=\tfrac{31}{100}$.  The sliver
$e\in[\tfrac{29}{100},\tfrac{31}{100}]$ (true ratio $\le0.13$) is closed by a
\textbf{fixed} $(e,t)$ cell table, $t:=d/\delta$, using
Lemma~\ref{zc9:lem:mvt} and the Case-II validity condition below.

\begin{lemma}\label{zc9:lem:certIIval}
The Case-II floor $mK_2fd$ is a valid lower bound for $C_m\psi/(\alpha^3p^{n})$
whenever $2(1+\rho)\xi\ge1$ (weaker than $2\rho\xi>1$).
\end{lemma}
\begin{proof}
The $\lambda=1$ dual (cert-II, $\RII$) gives $C_m\psi\ge\sqrt{2\alpha}B_mf
(d-\tfrac{e^2}{16\alpha^2(1+\rho)})$; and
$d-\tfrac{e^2}{16\alpha^2(1+\rho)}\ge\tfrac d2\iff\xi\ge\tfrac1{2(1+\rho)}$; then
relax $\tfrac12\sqrt{2\alpha}B_mf\ge\alpha^3mK_2f$ as in $\ZC$
Lem.~\texttt{payfloor}.
\end{proof}

\begin{proposition}[sliver]\label{zc9:prop:sliver}
\eqref{eq:target} holds for $m=9$ at every admissible point with
$\tfrac{29}{100}\le e\le\tfrac{31}{100}$ and $d<\delta$.
\end{proposition}

\begin{proof}
First, $R_9\le0$ on $\{t\le\tfrac{16}{25}\}$ throughout
$e\in[\tfrac{29}{100},\tfrac{31}{100}]$ (Part~6: grid max $-1.4\cdot10^{-6}$); on
that region \eqref{eq:target} is trivial.  On $t\in[\tfrac{16}{25},1]$ use the
fixed $4\times4$ cell table of Table~\ref{zc9:tab:sliver} ($e$-strips
$\tfrac{29}{100},\tfrac{59}{200},\tfrac3{10},\tfrac{61}{200},\tfrac{31}{100}$;
$t$-bands $\tfrac{16}{25},\tfrac{73}{100},\tfrac{41}{50},\tfrac{91}{100},1$).  On
each cell $\Def_9$ is bounded above by
$\tfrac{n x_U^{\,n-1}}{(\alpha p)^-}d_{\mathrm{hi}}(1-\ell_L^{\,m-1})$
(Lemma~\ref{zc9:lem:mvt} with $N\le d(1-\ell^{m-1})$, valid since $2d\le f$ by
Lemma~\ref{zc9:lem:cornerN}(i)); the payment is bounded below by
$mc_e d_{\mathrm{lo}}^2/e^2$ where $2\rho\xi\le1$ and by $mK_2fd_{\mathrm{lo}}$
where $2(1+\rho)\xi\ge1$ (Lemma~\ref{zc9:lem:certIIval}), the two conditions
covering each cell.  Every tabulated ratio is $<1$ (worst $0.690$), so
$\Def_9\le$ payment on each cell.
\end{proof}

\begin{table}[h]\centering\small
\caption{$m=9$ sliver: $16$ floored $(e,t)$ cells on
$[\tfrac{29}{100},\tfrac{31}{100}]\times[\tfrac{16}{25},1]$; $R_9\le0$ below
$t=\tfrac{16}{25}$.  Ratio $=\Def_9/\text{floor}$ (rounded up).}
\label{zc9:tab:sliver}
\begin{tabular}{lllll}\toprule
$e_L$&$e_R$&$t_L$&$t_R$&ratio $\le$\\\midrule
$\tfrac{29}{100}$&$\tfrac{59}{200}$&$\tfrac{16}{25}$&$\tfrac{73}{100}$&$0.362$\\
$\tfrac{29}{100}$&$\tfrac{59}{200}$&$\tfrac{73}{100}$&$\tfrac{41}{50}$&$0.353$\\
$\tfrac{29}{100}$&$\tfrac{59}{200}$&$\tfrac{41}{50}$&$\tfrac{91}{100}$&$0.346$\\
$\tfrac{29}{100}$&$\tfrac{59}{200}$&$\tfrac{91}{100}$&$1$&$0.340$\\
$\tfrac{59}{200}$&$\tfrac3{10}$&$\tfrac{16}{25}$&$\tfrac{73}{100}$&$0.690$\\
$\tfrac{59}{200}$&$\tfrac3{10}$&$\tfrac{73}{100}$&$\tfrac{41}{50}$&$0.605$\\
$\tfrac{59}{200}$&$\tfrac3{10}$&$\tfrac{41}{50}$&$\tfrac{91}{100}$&$0.343$\\
$\tfrac{59}{200}$&$\tfrac3{10}$&$\tfrac{91}{100}$&$1$&$0.337$\\
$\tfrac3{10}$&$\tfrac{61}{200}$&$\tfrac{16}{25}$&$\tfrac{73}{100}$&$0.000$\\
$\tfrac3{10}$&$\tfrac{61}{200}$&$\tfrac{73}{100}$&$\tfrac{41}{50}$&$0.638$\\
$\tfrac3{10}$&$\tfrac{61}{200}$&$\tfrac{41}{50}$&$\tfrac{91}{100}$&$0.343$\\
$\tfrac3{10}$&$\tfrac{61}{200}$&$\tfrac{91}{100}$&$1$&$0.338$\\
$\tfrac{61}{200}$&$\tfrac{31}{100}$&$\tfrac{16}{25}$&$\tfrac{73}{100}$&$0.000$\\
$\tfrac{61}{200}$&$\tfrac{31}{100}$&$\tfrac{73}{100}$&$\tfrac{41}{50}$&$0.685$\\
$\tfrac{61}{200}$&$\tfrac{31}{100}$&$\tfrac{41}{50}$&$\tfrac{91}{100}$&$0.613$\\
$\tfrac{61}{200}$&$\tfrac{31}{100}$&$\tfrac{91}{100}$&$1$&$0.344$\\
\bottomrule\end{tabular}
\end{table}

\subsection{The corner $e^\star_m\le e<\tfrac13$}\label{zc9:sec:wide}

Here $\delta=\tfrac{1-3e}6\in(0,\Delta_m]$, $\alpha=\tfrac13+\delta$,
$d\in(0,\delta)$ (no use of $\xi\le1$).  We use Lemma~\ref{zc9:lem:mvt} together
with the near-cancellation $d\approx q-L$ forced by $R_m>0$.

\begin{lemma}[corner defect]\label{zc9:lem:cornerN}
For any admissible point with $0<d<\delta$ (i.e.\ $q>\tfrac13$), $n=m-2$:
\begin{enumerate}
\item[(i)] $2d\le f$, equivalently $d\le q-L$.
\item[(ii)] If $R_m>0$ then $0<N\le d(1-\ell^{\,m-1})$, $N$ as in
 Lemma~\ref{zc9:lem:mvt}.
\item[(iii)] $\alpha\Def_m\le\tfrac{n x^{n-1}}p N$.
\end{enumerate}
\end{lemma}

\begin{proof}
(i) $2d\le f=\alpha-L\iff L\le2q-\alpha$; since $2q-\alpha\ge\tfrac13-\delta>0$,
this is $(2q-\alpha)^2-L^2\ge0$.  With $u:=q-\tfrac13=\delta-d>0$ and
$\alpha=\tfrac13+\delta$, direct expansion gives
$5q^2-q-4q\alpha+2\alpha^2=u+(5u^2-4u\delta+2\delta^2)>0$ (discriminant
$16\delta^2-40\delta^2<0$).  Then $d\le q-L$ from $2d\le\alpha-L=(q-L)+d$.
(ii) $R_m>0\Rightarrow N>0$; and $d\le q-L$ gives
$N=d-\ell^{m-1}(q-L)\le d(1-\ell^{m-1})$.  (iii) is Lemma~\ref{zc9:lem:mvt}.
\end{proof}

Let $\alpha(u)=\tfrac{1-u}2$ and, on $[e^\star_m,\tfrac13)$, put
\[
 a_U=\tfrac13+\Delta_m,\quad x_U=\tfrac{a_U}{1-a_U},\quad
 q_L^-=\tfrac23-2\Delta_m,\quad q_L^+=\tfrac23+\tfrac{\Delta_m}2,\quad
 (\alpha p)^-=\tfrac13(1-a_U),
\]
\[
 f_{\!H}=\tfrac{3a_U+\frac13}{q_L^-},\quad E_L=3f_{\!H},\quad
 \ell_L=1-E_L\Delta_m,\quad Q_L=\tfrac{2/3}{q_L^+},\quad
 c^-=\tfrac{2(1-x_U^{m-1})(1-2^{-(m-1)})}{(1+x_U)(3/2)},\quad
 K^-=\tfrac{\sqrt{2/3}\,(1-2^{-(m-1)})}{2a_U^3(3/2)}
\]
(all monotone-endpoint, checked in Part~5).

\begin{proposition}[corner]\label{zc9:prop:wide}
For each $m\in\{9,11,13\}$, \eqref{eq:target} holds at every admissible point
with $e^\star_m\le e<\tfrac13$ (indeed all $d<\delta$).
\end{proposition}

\begin{proof}
Assume $R_m>0$.  By Lemma~\ref{zc9:lem:cornerN}(ii)--(iii),
$\Def_m\le\tfrac{n x^{n-1}}{\alpha p}d(1-\ell^{m-1})$.
\emph{Case I} ($2\rho\xi\le1$): floor $mc_e\kappa^2\ge9mc_e d^2$
($\kappa\ge3d$).  Using $d\ge\ell^{m-1}(q-L)\ge\ell_L^{m-1}Q_L\delta$,
$1-\ell^{m-1}\le(m-1)(1-\ell)\le(m-1)E_L\delta$, $e^2\le\tfrac19$,
\[
 \frac{\Def_m}{mc_e\kappa^2}\ \le\
 \frac{n\,x_U^{\,n-1}}{(\alpha p)^-}\cdot\frac{(m-1)E_L}
 {9\,m\,c^-\,\ell_L^{\,m-1}\,Q_L}\ =:\ r_{\mathrm I}(m).
\]
\emph{Case II} ($2\rho\xi>1$): floor $mK_2fd$; with $f=\alpha(1-\ell)$,
$(1-\ell^{m-1})/f\le(m-1)/\alpha\le3(m-1)$,
\[
 \frac{\Def_m}{mK_2fd}\ \le\
 \frac{n\,x_U^{\,n-1}}{(\alpha p)^-}\cdot\frac{3(m-1)}{m\,K^-}\ =:\
 r_{\mathrm{II}}(m).
\]
Numerically (Part~5, exact envelopes):
\[
\begin{array}{lccccc}
\toprule
m & \Delta_m & x_U & \ell_L & r_{\mathrm I}(m) & r_{\mathrm{II}}(m)\\
\midrule
13 & \tfrac{13}{600} & 0.5504 & 0.8542 & 0.7063 & 0.0594\\
11 & \tfrac{9}{600}  & 0.5345 & 0.9026 & 0.5893 & 0.1171\\
 9 & \tfrac{7}{600}  & 0.5267 & 0.9256 & 0.9345 & 0.2765\\
\bottomrule
\end{array}
\]
Both are $<1$, so \eqref{eq:target} holds; if $R_m\le0$ it is trivial.
\end{proof}

\begin{proof}[Proof of Theorem~\ref{zc9:thm:main}]
Fix $m$ and an admissible point with $\xi\le1$, $q>\tfrac13$.  If
$e\le\tfrac1{60}$: Corollary~\ref{zc9:cor:za}.  If
$\tfrac1{60}\le e\le\bar e_m$: Proposition~\ref{zc9:prop:strip}.  If
$e\ge e^\star_m$: Proposition~\ref{zc9:prop:wide} (for $m\in\{11,13\}$,
$\bar e_m=e^\star_m$).  For $m=9$ the sliver
$[\tfrac{29}{100},\tfrac{31}{100}]$ is Proposition~\ref{zc9:prop:sliver} and
$[\tfrac{31}{100},\tfrac13)$ is Proposition~\ref{zc9:prop:wide}.  These cover
$(0,\tfrac13)$.
\end{proof}

\section{Assembly: the three zones tile the domain}\label{sec:assembly}

\begin{lemma}[the strip is empty beyond $e_1$]\label{lem:e1}
For $e>e_1=\tfrac{2033}{10000}$ there is no admissible $\kappa$ with $\xi\ge1$:
$\kxi(e)=\tfrac{e}{(1-e)^2}>\khi(e)=\min(\kmx(e),\kq(e))$.  Hence every admissible
point with $\xi\ge1$ has $e\le e_1$.
\end{lemma}

\begin{proof}
$\xi\ge1\iff\kappa\ge\kxi(e)$, and admissibility forces $\kappa\le\khi(e)$
(indeed $\kappa\le\kmx$ and $\kappa<\kq$).  So $\xi\ge1$ requires
$\kxi(e)\le\khi(e)$.  On $(0,\tfrac13)$, $\kxi$ is increasing and $\khi$
decreasing, and at $e=e_1$ one checks the exact rational comparison
$\kxi(e_1)>\khi(e_1)$ ($\khi(e_1)=\kq(e_1)=\tfrac{1-3e_1}{6e_1}$; Part~8 of the
script verifies $\kxi(e)>\khi(e)$ for all $e\in[e_1,\tfrac13)$).  Thus for
$e>e_1$ the strip is empty.
\end{proof}

\begin{proof}[Proof of Theorem~\ref{thm:main}]
Fix $m\in\{9,11,13\}$ and an admissible $(q,\alpha)$, i.e.\ $e\in(0,\tfrac13)$,
$0<\kappa\le\kmx(e)$, $\kappa<\kq(e)$.  If $R_m\le0$ then \eqref{eq:target} is
trivial ($C_m,\psi\ge0$).  Assume $R_m>0$ and split on $\xi=(1-e)^2\kappa/e$:

\smallskip\noindent\emph{(i) $\xi\le1$.}  Theorem~\ref{zc9:thm:main} (Zone~C)
gives \eqref{eq:target} for every admissible point with $\xi\le1$ and
$q>\tfrac13$ (and $q>\tfrac13$ is $\kappa<\kq$, part of admissibility).

\smallskip\noindent\emph{(ii) $\xi\ge1$ and $e\le\tfrac1{60}$.}
Proposition~\ref{prop:smalle} (Zone~A) gives \eqref{eq:target}.

\smallskip\noindent\emph{(iii) $\xi\ge1$ and $e\ge\tfrac1{60}$.}  By
Lemma~\ref{lem:e1}, $e\le e_1$, and $\xi\ge1$ is $\kappa\ge\kxi(e)$, while
admissibility gives $\kappa\le\khi(e)$.  So $(e,\kappa)$ is a strip point of
Proposition~\ref{zb9:battle} (Zone~B), which gives \eqref{eq:target}.

\smallskip The three cases exhaust the admissible domain: (i) covers $\xi\le1$;
(ii)--(iii) cover $\xi\ge1$ (every $\xi\ge1$ point has $e\le\tfrac1{60}$ or
$e\ge\tfrac1{60}$), and the boundaries $\xi=1$, $e=\tfrac1{60}$ lie in both
adjacent zones.  Finally Remark~\ref{rem:muniform} converts \eqref{eq:target}
into $t(C_m,W)\ge p^m-pq^{m-1}$ in Region~II.
\end{proof}

\begin{remark}[the zone target inequalities coincide]\label{rem:coincide}
The displayed targets \eqref{zb9:reduced} (Zone~B, via Lemma~\ref{zb9:red}) and
the Zone~C payment-floor inequalities are, in each case, sufficient conditions
for the \emph{same} inequality \eqref{eq:target}; the three zone statements
(Propositions~\ref{prop:smalle},~\ref{zb9:battle} and
Theorem~\ref{zc9:thm:main}) are the same conclusion \eqref{eq:target} on
complementary sub-domains.
\end{remark}

\section{Provenance: what is analytic and what remains}\label{sec:prov}

The result \eqref{eq:target} for $m\in\{9,11,13\}$ is established with no adaptive
or greedy interval subdivision anywhere.  Precisely:
\begin{itemize}
\item \textbf{Reduction, gate, Zone~A} (\S\ref{sec:relax}--\ref{sec:zoneA}):
 fully self-contained prose, following the small-$m$ note
 \texttt{smallm\_regionII\_proof\_v2.tex} verbatim; the only numerics are the
 exact-rational constants of Table~\ref{tab:smalle} (all safe-rounded).
\item \textbf{Zone~B} (\S\ref{sec:zoneB}): a $26$-row displayed certificate
 (Tables~\ref{zb9:tab9}--\ref{zb9:tab13}) over a \emph{fixed} rational cover;
 each row is a closed-form monotone-endpoint comparison $L(E)\ge R(E)$.  This
 meets the project's analytic bar (a displayed fixed-cover table).  It quotes,
 by name, the monotonicity toolkit and reduction of $\ZB$/$\RII$; those are the
 verified analytic $m\ge15$ templates, here specialised (no $m$-maximisation).
\item \textbf{Zone~C strip and corner} (\S\ref{zc9:sec:strip},
 \S\ref{zc9:sec:wide}): the strip is a fixed $\varphi$-row table
 (Tables~\ref{zc9:tab:m13}--\ref{zc9:tab:m9}, $66$ rows) over the $\ZC$ cover
 plus displayed bisections; the corner is a $1$-D dual-certificate chain
 $r_{\mathrm I},r_{\mathrm{II}}<1$ in closed form.  The single ingredient beyond
 $\ZC$ is the mean-value defect (Lemma~\ref{zc9:lem:mvt}), proved here in full.
 The strip row bounds $\Psi_{\mathrm{deep}},\Psi_{\mathrm{sh}},\Psi_{\mathrm{II}}$
 are elementary (exp/log/sqrt of rational endpoints), evaluated once per fixed
 interval and safe-rounded; they are not reduced to bare integer comparisons.
\item \textbf{Zone~C $m=9$ sliver} (\S\ref{zc9:sec:sliver}): a \emph{fixed}
 $16$-cell $(e,t)$ table (Table~\ref{zc9:tab:sliver}).  It is a fixed cover (not
 adaptive), but it is the least self-contained piece: the per-cell
 monotone-envelope constants $d_{\mathrm{hi}},d_{\mathrm{lo}},\ell_L,x_U$ and the
 Case-I/Case-II branch selection are generated and checked by the script rather
 than written out symbolically in the note.  The displayed ratios ($<1$, worst
 $0.690$) are what the certificate asserts.
\end{itemize}
The two source certificates were each independently machine-verified
(\texttt{validate\_zoneB\_91113.py}, \texttt{validate\_zoneC\_91113\_v2.py});
the merged script \texttt{validate\_m91113.py} re-runs both and additionally
verifies the tiling (Lemma~\ref{lem:e1}) and the end-to-end bound
$R_m\le C_m\psi$ over the \emph{entire} admissible domain (dense
$(e,\kappa)$ grid): the observed grid maximum of $R_m/(C_m\psi)$ is
$\approx0.71,\,0.76,\,0.80$ for $m=9,11,13$ (attained near the small-$e$,
high-$\xi$ pinch), comfortably below $1$; the maximum over the $\xi\ge1$
strip alone is $0.583,\,0.573,\,0.554$.  No step depends on an unverified
numeric.

\emph{Residues.}  (1) The note is a \emph{specialisation}: it cites, without
re-proving, named lemmas of $\RII$, $\ZB$, $\ZC$ (the verified analytic
$m\ge15$ templates).  (2) The $m=9$ sliver relies on a fixed but
script-generated $2$-D cell table, heavier than a written-out $1$-D
$\varphi$-row.  Neither is an adaptive subdivision; both meet the fixed-cover
bar.

\end{document}
